\begin{textsample}{POS Dim 1 – 2013\_c4 – Score -5.00 – 2013\_c4/enade\_2013\_enf\_7b93b4bd5c.txt}  \label{ex:f1_pos_389}
CAPÍTULO 4 PERCEPÇÃO DA PROVA As análises feitas neste capítulo tratam das percepções dos concluintes da Área de Enfermagem sobre a prova aplicada no ENADE/2013.
Estas percepções foram mensuradas por meio de nove questões que avaliaram desde o grau de dificuldade da prova até o tempo gasto para concluí -la.
As percepções sobre a prova foram relacionadas com o desempenho dos estudantes e com a Grande Região de funcionamento do curso.
O questionário de percepção da prova encontra -se ao final do Anexo IV, que traz a reprodução da prova.
O desempenho dos estudantes foi classificado em quatro quartos.
Para tanto, esse desempenho foi ordenado de forma ascendente.
O percentil 25, P25, também conhecido como primeiro quartil, é a nota de desempenho que deixa um quarto ( 25 \% ) dos valores observados abaixo e três quartos acima.
A Figura 4.1 apresenta uma ilustração deste conceito.
O quarto inferior de desempenho é composto pelas notas abaixo do primeiro quartil.
Já o percentil 75, P75, também conhecido como terceiro quartil, é o valor para o qual há três quartos ( 75 \% ) dos dados abaixo e um quarto acima dele.
O quarto superior de desempenho é composto pelas notas iguais ou acima do terceiro quartil.
O percentil 50, P50, também conhecido como mediana, é o valor que divide as notas em dois conjuntos de igual tamanho.
O segundo quarto inclui valores entre o primeiro quartil ( P25 ) e a mediana.
O terceiro quarto contém os valores entre a mediana ( P50 ) e o terceiro quartil ( P75 ).
Vale ressaltar que percentis, quartis e medianas são pontos que não obrigatoriamente pertencem ao conjunto original de dados, ao passo que os quartos são subconjuntos dos dados originais.
A seguir, serão apresentados gráficos com resultados selecionados, relativos às nove questões avaliadas por grupos de estudantes.
Os gráficos apresentam nas barras o percentual de alunos que assinalaram uma das opções ou a soma das porcentagens daqueles que assinalaram duas ( ou três ) delas.
Por exemplo, para as questões 1 e 2, os gráficos apresentam a porcentagem total de participantes que assinalaram as opções ( D ) difícil e ( E ) muito difícil.
Em cada barra foram assinalados também os extremos do intervalo de confiança de 95 \% como linhas verticais unidas por uma linha horizontal na forma da letra H maiúscula, semelhantemente aos gráficos do Capítulo 3.
As Tabelas no Anexo II apresentam os valores absolutos e a distribuição percentual12 das alternativas válidas das nove questões, segundo o mesmo recorte de desempenho dos alunos e Grande Região de funcionamento do curso. 4.1 GRAU DE DIFICULDADE DA PROVA 4.1.1 Componente de Formação Geral Ao avaliarem “ Qual o grau de dificuldade desta prova na Parte de Formação Geral? ” ( Questão 1 ), 27,6 \% do grupo de inscritos e presentes optaram pelas alternativas difícil ou muito difícil.
Entretanto, para mais da metade dos estudantes ( 63,3 \% ), o Componente de Formação Geral da prova foi considerado com grau de dificuldade médio ( Gráfico 4.1, Gráfico 4.2 e, no Anexo II, a Tabela II.1 ).
O percentual de estudantes que consideraram a prova como difícil ou muito difícil foi maior na região Sul, onde a proporção foi de 32,1 \%, enquanto a de menor incidência foi a Nordeste, com 22,8 \%.
No Gráfico 4.1 é possível observar que esta diferença é estatisticamente significativa.
A diferença entre a proporção de respondentes com o agregado destas opções na região Nordeste é estatisticamente diferentes das proporções observadas nas demais regiões.
Nas Grandes Regiões, a proporção de presentes à prova que consideraram o Componente de Formação Geral como sendo de grau de dificuldade médio esteve entre 60,0 \% ( região Sul ) e 66,6 \% ( região Nordeste ), situação inversa dos que consideraram a prova como difícil ou muito difícil. 12 Cumpre lembrar uma das convenções para tabelas numéricas ( pág. iii ) sobre a possibilidade da soma das partes não resultar em 100 \% por questões de arredondamento.
O percentual de alunos que consideraram a prova difícil ou muito difícil foi decrescente em relação ao aumento de desempenho ( Gráfico 4.2 ).
Nos dois quartos de menor desempenho, com diferença estatisticamente significativa, 34,2 \% no 1º quarto e 29,5 \% no 2º consideraram a prova difícil ou muito difícil.
Nos quartos de maior desempenho a proporção de alunos que julgaram a prova difícil ou muito difícil – 26,7 \% no 3º quarto e 20,8 \% no 4º quarto – foi menor do que nos demais e significativamente diferente entre si.
A diferença entre o 2º e o 3º quarto também foi estatisticamente significativa.
Para todos os quartos de desempenho a alternativa modal para esta pergunta foi médio, com 59,1 \%, 63,0 \%, 64,4 \% e 66,5 \% dos respondentes de cada um dos quartos, de 1 a 4 respectivamente, com valores crescentes com o desempenho. 4.1.2 Componente de Conhecimento Específico Ao responderem à Questão 2 – “ Qual o grau de dificuldade desta prova na parte de Componente Específico? ” – 28,1 \% do grupo de estudantes classificaram -na como difícil ou muito difícil.
Além disso, o Componente de Conhecimento Específico da prova foi considerado com grau de dificuldade médio por 64,2 \% dos alunos ( Gráfico 4.3, Gráfico 4.4 e, no Anexo II, a Tabela II.2 ).
A análise das respostas dos estudantes quanto ao grau de dificuldade do Componente de Conhecimento Específico da prova, agregados por Grande Região, mostra que a diferença entre a maior e a menor proporção de alunos que a avaliaram como difícil ou muito difícil é estatisticamente significativa : 24,4 \% na região Nordeste e 31,3 \% nas regiões Sul e Centro-Oeste.
O percentual de alunos que classificaram o grau de dificuldade como médio, no Componente de Conhecimento Específico, variou de 61,5 \% a 67,1 \%, para as regiões Centro-Oeste e Nordeste, respectivamente, novamente situação inversa dos que consideraram este componente da prova como difícil ou muito difícil.
Considerando -se a avaliação da dificuldade das questões do Componente de Conhecimento Específico da prova, de acordo com o desempenho dos estudantes, observam -se diferenças estatisticamente significativas entre alguns quartos, como a existente entre o quarto superior e os demais.
A proporção dos que classificaram a parte específica como difícil ou muito difícil variou de 31,4 \% ( 1º quarto ) a 24,2 \% ( 4º quarto ), decrescente com o desempenho.
Já a alternativa modal para a Questão 2 foi o grau médio, com 61,5 \% do quarto inferior e 66,5 \% do superior optando por esta resposta, apresentando tendência de crescimento conforme o desempenho. 4.2 EXTENSÃO DA PROVA EM RELAÇÃO AO TEMPO TOTAL Indagados quanto à extensão da prova, em relação ao tempo total oferecido para a sua resolução ( Questão 3 ), os estudantes apontaram, com maior incidência, as alternativas que considerava a extensão longa ou muito longa ( Gráfico 4.5, Gráfico 4.6 e, no Anexo II, a Tabela II.3 ).
O percentual de alunos que responderam ser a extensão da prova adequada foi de 44,7 \%.
Já 48,6 \% dos inscritos presentes consideraram que o exame foi longo ou muito longo e 6,7 \% o avaliaram como curto ou muito curto.
Entre as Grandes Regiões a proporção daqueles que avaliaram a prova como longa ou muito longa em relação ao tempo total destinado à sua resolução variou de 38,4 \% na região Norte até 51,1 \% na região Sudeste.
A proporção dos que consideraram a prova longa ou muito longa na região Norte é significativamente diferente da proporção dos que deram as mesmas respostas nas demais regiões, o mesmo ocorrendo com a região Centro-Oeste.
Considerando -se o desempenho dos alunos, nota -se ainda que o percentual dos que consideraram a extensão da prova adequada variou de 46,6 \%, no 4º quarto de desempenho, a 42,7 \% no quarto de desempenho inferior.
No Gráfico 4.6, pode -se constatar que para os quatro quartos de desempenho a proporção de estudantes que consideraram a prova longa ou muito longa, em relação ao tempo total destinado à sua resolução, manteve -se em torno de 48 \%, sem diferenças estatisticamente significativas, mas apresentando, grosso modo, uma tendência decrescente com o desempenho. 4.3 COMPREENSÃO DOS ENUNCIADOS DAS QUESTÕES 4.3.1 Componente de Formação Geral Com relação aos enunciados das questões do Componente de Formação Geral ( Questão 4 ), as opiniões foram positivas, já que 72,6 \% dos alunos avaliados consideraram os enunciados de todas ou da maioria das questões claros e objetivos ( Gráfico 4.7, Gráfico 4.8 e, no Anexo II, a Tabela II.4 ).
Na análise regional, a percentagem de estudantes que avaliaram que todos ou a maioria dos enunciados das questões do Componente de Formação Geral estavam claros e objetivos variou de 69,6 \% na região Sudeste a 76,1 \% na região Centro-Oeste, sendo esta diferença estatisticamente significativa.
A análise das percepções dos estudantes sobre a clareza e objetividade dos enunciados permite afirmar que todos ou a maioria dos enunciados de questões relativas ao Componente de Formação Geral foram considerados claros e objetivos para a maior parte dos respondentes ( maior do que 69 \% em todas as regiões e maior do que 64 \% para todos os quartos de desempenho ).
Segundo o desempenho, observa -se que a proporção dos que emitiram esta opinião cresce conforme o desempenho aumenta, com diferenças estatisticamente significativas entre todos os quartos de desempenho.
No quarto superior, a clareza e objetividade de todos ou da maioria dos enunciados das questões foi percebida por 80,3 \% dos alunos e no quarto de desempenho inferior tal avaliação foi emitida por 64,4 \% deles. 4.3.2 Componente de Conhecimento Específico Com relação aos enunciados das questões do Componente de Conhecimento Específico da prova, para 79,2 \% dos estudantes avaliados da Área de Enfermagem a clareza e a objetividade ( Questão 5 ) estavam presentes em todas ou na maioria das questões ( Gráfico 4.9, Gráfico 4.10 e, no Anexo II, a Tabela II.5 ).
A maioria dos estudantes de todas as Grandes Regiões brasileiras considerou claros e objetivos todos ou a maioria dos enunciados das questões do Componente de Conhecimento Específico da prova, percentual sempre maior do que 76 \%.
As diferenças entre as regiões Norte, Nordeste e Centro-Oeste não são estatisticamente significativas.
As diferenças dessas proporções entre as regiões Sul ( 82,8 \%, a maior ) e Sudeste ( 76,7 \%, a menor ) e as demais são todas estatisticamente significativas.
A proporção de estudantes que consideraram os enunciados das questões claros e objetivos apresenta uma tendência crescente em relação ao aumento de desempenho : mais elevada no quarto superior ( 87,8 \% ) se comparada ao quarto inferior de desempenho ( 69,2 \% ).
As diferenças dos resultados entre os quartos de desempenho são estatisticamente significativas. 4.4 SUFICIÊNCIA DAS INFORMAÇÕES/INSTRUÇÕES FORNECIDAS Ao avaliarem as informações/instruções fornecidas para a resolução das questões ( Questão 6 ), 84,2 \% dos respondentes da Área de Enfermagem de todo o Brasil afirmaram que estas eram até excessivas ou suficientes em todas ou na maioria das questões ( Gráfico 4.11, Gráfico 4.12 e, no Anexo II, a Tabela II.6 ).
Quanto à distribuição de respondentes pelas Grandes Regiões observa -se que a proporção de estudantes que consideraram as informações/instruções fornecidas até excessivas ou suficientes em todas ou na maioria das questões foi sempre superior a 82 \%, chegando a 87,0 \% na região Sul.
As diferenças entre as proporções das regiões Norte ( 82,5 \% ) e Sudeste ( 82,7 \% ), as menores proporções, e a da região Sul, a maior proporção, são estatisticamente significativas.
Levando -se em conta o desempenho dos participantes, notam -se diferenças estatisticamente significativas das opiniões de estudantes entre todos os quartos de desempenho, como mostra o Gráfico 4.12.
O percentual de participantes que avaliaram as informações/instruções como até excessivas ou suficientes em todas ou na maioria das questões foi mais elevado no quarto superior ( 91,5 \% ), percentual superior à média nacional ( 84,2 \% ).
Já no quarto inferior, a suficiência das informações/instruções declarada como até excessiva em todas ou na maioria das questões foi percebida por 75,1 \% dos respondentes, com um comportamento crescente com os quartos de desempenho. 4.5 DIFICULDADE ENCONTRADA AO RESPONDER À PROVA Perguntados sobre as dificuldades com as quais se depararam ao responder à prova ( Questão 7 ), 6,1 \% dos estudantes apontaram o desconhecimento do conteúdo.
Para 47,0 \%, a forma diferente de abordagem do conteúdo foi indicada como dificuldade.
Já a falta de motivação para fazer a prova foi a dificuldade apontada por 16,8 \% dos respondentes.
Considerando -se todo o Brasil, 25,5 \% dos respondentes afirmaram que não tiveram qualquer tipo de dificuldade para responder à prova ( Tabela II.7 no Anexo II ).
Os Gráficos 4.13 e 4.14 apresentam os percentuais de estudantes que apontaram o desconhecimento do conteúdo como dificuldade percebida ao responder à prova.
Na análise por Grandes Regiões, o percentual de inscritos e presentes que apontaram o desconhecimento do conteúdo como dificuldade ao responder à prova não superou 9 \%.
Os percentuais variaram de 5,0 \% na região Nordeste a 8,4 \% no Sul, e esta diferença é estatisticamente significativa.
A forma diferente de abordagem do conteúdo foi a escolha modal dos estudantes, com percentuais que variaram de 43,6 \% ( região Nordeste ) a 53,0 \% ( região Norte ).
O percentual de alunos que citaram a falta de motivação como dificuldade variou de 13,1 \% ( região Norte ) a 18,4 \% ( região Sudeste ).
A proporção dos que declararam não ter qualquer dificuldade para responder à prova variou de 22,5 \% na região Norte a 29,0 \% na região Nordeste.
Com relação aos quartos de desempenho, o desconhecimento do conteúdo, foi a opção escolhida por 6,6 \% dos estudantes do quarto superior e 6,2 \% do quarto inferior.
A alternativa modal para os alunos, quando agregados pelos quartos de desempenho, foi que a dificuldade encontrada foi causada pela forma diferente de abordagem do conteúdo, com valores entre 42,4 \% no quarto superior e 49,8 \% no 2º quarto. 4.6 CONTEÚDOS DAS QUESTÕES OBJETIVAS DA PROVA Ao analisarem os conteúdos das questões objetivas da prova ( Questão 8 ), um percentual muito pequeno dos estudantes avaliados, apenas 3,3 \%, afirmou que não estudou ainda a maioria desses conteúdos ( Gráfico 4.15, Gráfico 4.16 e a Tabela II.8 no Anexo II ).
A maioria ( 77,5 \% ) afirmou ter estudado e aprendido muitos ou todos os conteúdos.
Na análise por Grande Região, a proporção de respondentes que escolheram a opção não estudou ainda a maioria desses conteúdos, foi também pequena.
Observa -se que nas regiões Norte ( 5,2 \% ) e Sudeste ( 3,4 \% ), apesar de pequenas, as proporções foram maiores do que a média nacional ( 3,3 \% ).
As regiões Nordeste ( 3,2 \% ), Sudeste ( 3,4 \% ) e Centro-Oeste ( 3,3 \% ) têm proporções muito próximas ou idêntica à média nacional.
Observam -se diferença estatisticamente significativa entre a região Norte e as demais regiões.
Em todas as regiões, a maioria dos presentes afirmou ter estudado e aprendido muitos ou todos os conteúdos, com proporções variando entre 71,7 \% na região Norte e 79,9 \% na região Sul.
Considerando -se separadamente as opiniões de estudantes dos quatro quartos de desempenho, observa -se que, no quarto inferior, 7,0 \% ofereceram como resposta que não estudou ainda a maioria desses conteúdos, sendo 1,1 \% os do quarto superior com a mesma resposta, com valores decrescentes com o desempenho.
As diferenças entre os alunos que optaram por este motivo de dificuldade são estatisticamente significativas entre todos os quartos de desempenho.
Tendo em conta o quarto superior, 88,3 \% dos alunos afirmaram ter estudado e aprendido muitos ou todos os conteúdos.
No outro extremo, no quarto inferior, 63,2 \% dos alunos responderam da mesma forma, com valores crescentes com o desempenho. 4.7 TEMPO GASTO PARA CONCLUIR A PROVA Ao responderem sobre o tempo de conclusão da prova ( Questão 9 ), 3/4 dos estudantes ( 75,0 \% ) afirmou ter gasto entre duas e quatro horas ( Gráfico 4.17, Gráfico 4.18 e, no Anexo II, a Tabela II.9 ).
Considerando -se as cinco Grandes Regiões brasileiras, nas regiões Norte ( 73,5 \% ) e Nordeste ( 74,2 \% ) o percentual dos que utilizaram entre duas e quatro horas para finalizar a prova foi inferior à média nacional ( 75,0 \% ).
Em duas das cinco Grandes Regiões, Sul e Centro-Oeste, o percentual de alunos que dispensaram entre duas e quatro horas para concluir a prova, superou 76 \%, como mostra o Gráfico 4.17.
Uma vez sendo analisadas as alternativas escolhidas pelos estudantes que se situam nos diferentes quartos de desempenho, observa -se que uma proporção bem maior de participantes no quarto superior declarou ter gasto entre duas e quatro horas para concluir a prova quando comparados com os do quarto inferior, respectivamente 81,7 \% e 65,9 \%.
As diferenças entre os quatro quartos são estatisticamente significativas e evidenciam uma tendência crescente.

% matched lemmas: 
\end{textsample}
